% !PRINTED 179
\weberpage{179}
4. Die Invariante $j(\omega)$ gehört zur Gruppe $\mathfrak{L}$ (\S\ 53) und
daher der Satz:

\gesperrt{Jede Modulfunktion, die durch die beiden Vertauschungen
}
\[
(\omega,\ \omega+1), \qquad \left(\omega,\ -\frac{1}{\omega}\right)
\]
\gesperrt{ungeändert bleibt, ist eine rationale Funktion von $j(\omega)$.}

Außer diesen führen wir noch zwei andere Modulfunktionen
ein.

Aus der Definition der Funktionen $f(\omega)$, $f_1(\omega)$, $f_2(\omega)$ [\S\,34,
(1), (10)] ergibt sich, wenn man die beiden Gleichungen \S\,42, (3)
miteinander multipliziert:
\weq{2}{f(\omega) = \frac{\sqrt[6]{2}}{\sqrt[12]{\varkappa\varkappa'}}, \qquad f_1(\omega) = \sqrt[6]{2}\ \sqrt[12]{\frac{\varkappa'^2}{\varkappa}}, \qquad f_2(\omega) = \sqrt[6]{2}\ \sqrt[12]{\frac{\varkappa^2}{\varkappa'}},}
\weq{3}{\frac{f_1(\omega)}{f(\omega)} = \sqrt[4]{\varkappa'}, \qquad \frac{f_2(\omega)}{f(\omega)} = \sqrt[4]{\varkappa}.}

Auf Grund von \S\,46, (15), (16) definieren wir zwei Funktionen
 $\gamma_2(\omega)$, $\gamma_3(\omega)$:
\weq{4}{\begin{aligned}
\gamma_2(\omega) &= \sqrt[3]{j(\omega)} = \sqrt[3]{2^8}\ \frac{1 - \varkappa^2\varkappa'^2}{\sqrt[3]{\varkappa^4\varkappa'^4}},\\
\gamma_3(\omega) &= \sqrt{j(\omega) - 27\cdot 64} = \frac{8(2+\varkappa^2\varkappa'^2)(\varkappa'^2 - \varkappa^2)}{\varkappa^2\varkappa'^2},
\end{aligned}}
die nach (2) und (3) als eindeutige Funktionen von $\omega$ folgendermaßen
 darstellbar sind:
\weq{5}{\begin{aligned}
\gamma_2(\omega) &= \frac{f(\omega)^{24} - 16}{f(\omega)^8},\\
\gamma_3(\omega) &= \frac{[f(\omega)^{24} + 8][f_1(\omega)^8 - f_2(\omega)^8]}{f(\omega)^8},
\end{aligned}}
wofür man nach den Grundformeln für die $f$-Funktionen [\S\,34,
(11), (12)] auch setzen kann:
\weq{6}{\gamma_2(\omega) = f(\omega)^8 f_1(\omega)^8 + f(\omega)^8 f_2(\omega)^8 - f_1(\omega)^8 f_2(\omega)^8,}
\weq{7}{\gamma_3(\omega) = \tfrac{1}{2}[f(\omega)^8 + f_1(\omega)^8][f(\omega)^8 + f_2(\omega)^8][f_1(\omega)^8 - f_2(\omega)^8].}
Es sind also $x = f^8,\ -f_1^8,\ -f_2^8$ die Wurzeln der kubischen
Gleichung
\weq{8}{x^3 - \gamma_2 x - 16 = 0,}
und $4\gamma_3^2$ ist die Diskriminante dieser Gleichung.
